\scp{Apicals}{12-02}
Like both \tsc{\ili{Ambon}-Seram}
and \tsc{Seram-Tanimbar-Bomberai}, \ili{Banda} also merges PMP *d/*z.
However, in order to determine the significance of this merger,
it needs to be discussed in the context of the set of historical
apical consonants: *d, *z, *R, *l, *j.
Looking at the whole set shows different groupings of mergers
and different reflexes, ultimately making it difficult to place
\ili{Banda} with either \tsc{\ili{Ambon}-Seram} or \tsc{Seram-Tanimbar-Bomberai}
based on the apicals.
\trf{12-02} summarises the reflexes of the apicals in languages
of central Maluku along with \tsc{Timor-Babar} to the south of the \ili{Banda} Islands.

\begin{table}\tcp{PMP apicals in \ili{Banda} and other subgroups}{12-02}
	\begin{threeparttable}\st{6pt}
	\begin{tabular}{lllllll}\lsptoprule
PMP &{\hr}*d&{\hr}*z&{\hr}*R&{\hr}*l&{\hr}*j\\	\midrule
\tsc{Sula-Buru}&**d&**z&**h&**l{\cdr}&\mc{2}{l}{**y, **l, {\0}}\\
\tsc{\ili{Ambon}-Seram}&**r{\cor}&**r{\cor}&**R&**l{\cdr}&**l{\cdr}&\\
\hspace{5mm}\tsc{Nuclear \ili{Ambon}-Seram}&**r{\cor}&**r{\cor}&**l{\cdr}&**l{\cdr}&**l{\cdr}&\\
\ili{Banda}&\hp{**}r{\cor}&\hp{**}r{\cor}&\hp{**}r{\cor}&\hp{**}l{\cdr}&\hp{**}l{\cdr}&\\
\tsc{Seram-Tanimbar-Bomberai}
&**d{\cbl}&**d\su{†}{\cbl}&**r{\cor}&**l{\cdr}&**j&\\
\tsc{Timor-Babar}&**d&**z&**R&**l{\cdr}&**j\\
\hspace{5mm}\tsc{\ili{Teun}-\ili{Nila}-Serua}&**d&**z&**r{\cor}&**l{\cdr}&**r{\cor}\\
		\lspbottomrule
	\end{tabular}
		\begin{tablenotes}\tnsize
			\item[†] {While \tsc{Seram-Tanimbar-Bomberai} does have initial *d/*z > **d,
								the data are a little ambiguous for medial *d/*z > **d.
								See \srf{11-01-01} for discussion.}
		\end{tablenotes}
	\end{threeparttable}
\end{table}

At the highest level,
we could place \tsc{\ili{Ambon}-Seram},
\tsc{Seram-Tanimbar-Bomberai}, and \ili{Banda} together on the basis
of shared merger of *d/*z,
if the changes affecting *R are viewed as secondary developments.
In particular we note that \tsc{\ili{Ambon}-Seram} *d/*z > **r would
be identical to \ili{Banda} *d/*z > \it{r}.
This, combined with the merger of *n/*ñ/*ŋ > \it{n} indicates
that the best avenue to explore for subgrouping \ili{Banda} would be
to link it to \tsc{\ili{Ambon}-Seram}.

\ili{As} shown in \trf{12-02},
\ili{Banda} merges PMP *d/*z/*R > \it{r},
and *l/*j > \it{l},
reflecting a different combination from other subgroups.
\ili{Banda} examples of the merger *d/*z/*R > \it{r} are shown in examples \qf{12-03}--\qf{12-05}.
Data from neighbouring languages in other subgroups can be seen
in Tables \ref{tab:07-02}, \ref{tab:07-05}, \ref{tab:07-07}
and \srf{04-06-04} (\tsc{\ili{Teun}-\ili{Nila}-Serua}).

{\wg\st{6pt}\begin{xltabular}[l]{\linewidth}
{>{\hspace{-\tabcolsep}}l<{\hspace{-\tabcolsep}}
>{\hspace{-\tabcolsep}\enspace}ll
>{\raggedright\arraybackslash}X}
\lsptoprule
						&	PMP				&	\ili{Banda}			&	\ili{Banda} gloss		\\ \midrule	\hhline{~}
%\endfirsthead
						%&	PMP			&	\ili{Funai}			&	\ili{Banda} gloss					\\
\endhead
\tbx{12-03}	&	*\tbr{d}ahun	&	\tbr{r}ano	&	leaf	\\
	&	*\tbr{d}aRaq	&	\tbr{r}aro-	&	blood	\\
	&	*\tbr{d}aRəq	&	\tbr{r}aro	&	clay	\\
	&	*\tbr{d}uha	&	\tbr{r}uo	&	two	\\
	&	*\tbr{d}uRi	&	\tbr{r}iri	&	bone	\\
	&	*\tbr{d}iŋdiŋ	&	\tbr{r}indin	&	cold	\\
	&	*m-u\tbr{d}əhi	&	fal-mu\tbr{r}i	&	behind	\\	\midrule	\hhline{~}
\tbx{12-04}	&	*\tbr{z}auq	&	\tbr{r}au	&	far	\\
	&	*\tbr{z}alan	&	e\tbr{r}an	&	road, path	\\
	&	*qa\tbr{z}ay	&	a\tbr{r}a-	&	chin, jaw	\\
	&	*qu\tbr{z}an	&	u\tbr{r}en	&	rain	\\
	&	*haRə\tbr{z}an	&	e\tbr{r}en	&	ladder	\\	\midrule	\hhline{~}
\tbx{12-05}	&	*\tbr{R}umaq	&	\tbr{r}umo	&	house	\\
	&	*da\tbr{R}aq	&	ra\tbr{r}o-	&	blood	\\
	&	*da\tbr{R}əq	&	ra\tbr{r}o	&	clay	\\
	&	*di\tbr{R}i	&	meli\tbr{r}i	&	stand	\\
	&	*du\tbr{R}i	&	ri\tbr{r}i	&	bone	\\
	&	*pa\tbr{R}ih	&	a\tbr{r}i	&	stingray	\\
	&	*wa\tbr{R}əj	&	wa\tbr{r}o-t	&	vine, rope	\\
	&	*baqə\tbr{R}u	&	fer{\tl}fe\tbr{r}u-ŋo	&	new	\\
	&	*haba\tbr{R}at	&	fa\tbr{r}at	&	west monsoon	\\
	&	*timu\tbr{R}	&	timu\tbr{r}	&	\ili{east} monsoon	\\
	&	*laya\tbr{R}	&	lea\tbr{r}	&	sail	\\
	&	*ma-lapa\tbr{R}	&	mla\tbr{r}, mala\tbr{r}	&	hungry 	\\	\hhline{~}
	&									&													&	(inflected; possible double vowel /mlaar/)	\\
	&	*qatəlu\tbr{R}	&	tulu\tbr{r}u	&	egg	\\
	&	*wahiyə\tbr{R}	&	wa\tbr{r}	&	water	\\
	&	*niu\tbr{R} > 	&	nue\tbr{r}	&	coconut	\\ %\hhline
	%&	> **niwə\tbr{R} 	&		&		\\
	&	**ŋi\tbr{R}ut	&	ni\tbr{r}ut-o	&	mucous	\\
\lspbottomrule
\end{xltabular}\addtocounter{table}{-1}}

Examples \qf{12-06}
and \qf{12-07} show the \ili{Banda} merger of PMP *l/*j > l.

\noindent{\wg\st{6pt}\tblr{>{\hspace{-\tabcolsep}}l<{\hspace{-\tabcolsep}}
>{\hspace{-\tabcolsep}\enspace}
lll}
\lsptoprule
\tbx{12-06}		& PMP & \ili{Banda} & \ili{Banda} gloss \\ \midrule
	&	*\tbr{l}ayaR	&	\tbr{l}ear {\tl} \tbr{l}iar	&	sail	\\
	&	*\tbr{l}ima	&	\tbr{l}imo	&	five	\\
	&	*ma-\tbr{l}apaR	&	m\tbr{l}ar, ma\tbr{l}ar	&	hungry 	\\	\hhline{~}
	&									&													&	(inflected; possible double vowel /mlaar/)	\\
	&	*a\tbr{l}aq	&	-a\tbr{l}a	&	take	\\
	&	*bu\tbr{l}an	&	fu\tbr{l}an	&	moon	\\
	&	*ti\tbr{l}u	&	ti\tbr{l}u-	&	ear	\\
	&	*tə\tbr{l}u	&	te\tbr{l}u	&	three	\\
	&	*qatə\tbr{l}uR	&	tu\tbr{l}uru	&	egg	\\
	&	*qu\tbr{l}u	&	u\tbr{l}u-n	&	head	\\
	&	**ma\tbr{l}ip	&	moma\tbr{l}ik	&	laugh, smile	\\
	&	**ku\tbr{l}a	&	ku\tbr{l}a	&	banana	\\	\midrule
\end{tabular}\mbox{}\\}

\noindent{\wg\st{6pt}\tblr{>{\hspace{-\tabcolsep}}l<{\hspace{-\tabcolsep}}
>{\hspace{-\tabcolsep}\enspace}
lll}
\lsptoprule
\tbx{12-07}	& PMP 						& \ili{Banda} 				& \ili{Banda} gloss \\ \midrule
	&	*ŋa\tbr{j}an	&	na\tbr{l}a-n	&	name	\\
	&	*pa\tbr{j}ay		&	a\tbr{l}a			&	rice	\\
	&	*pi\tbr{j}a			&	i\tbr{l}o			&	how many, how much	\\
	&	*qapə\tbr{j}u		&	e\tbr{l}u-n		&	gall, bile	\\
	&	*qalə\tbr{j}aw	&	lea						&	sun, day 	\\ \hhline{~}
	&									&								& (LS *j > {\0} also in \tsc{\ili{Ambon}-Seram} and \tsc{Sula-Buru})	\\
\lspbottomrule
\end{tabular}\mbox{}\\}

While we could link \ili{Banda} with \tsc{\ili{Ambon}-Seram} on the basis
of the shared merger of *d/*z,
the different reflexes and combination of the mergers involving
PMP *R makes this difficult as \tsc{Nuclear \ili{Ambon}-Seram} merges
PMP *R/*l/*j > **l (see \crf{ch:AS}).
\ili{Banda} could be placed as a primary branch of \tsc{\ili{Ambon}-Seram}
coordinate with \tsc{Nuclear \ili{Ambon}-Seram}, Boano, and \tsc{Seti}.
But even in Boano
and \tsc{Seti} there is a partial (though not full) merger of *R with
*l/*j, which is different to \ili{Banda} which does not merge *R with *l/*j.

The different result of *d/*z > **d
and *R > \it{r} for \tsc{Seram-Tanimbar-Bomberai} in contrast
to *d/*z/*R > \it{r} in \ili{Banda} suggests a closer look is needed,
since the changes affecting these segments in \tsc{Seram-Tanimbar-Bomberai}
and/or \ili{Banda} could more easily reflect secondary developments.

Indeed, it is the shared merger of *d/*z/*R > \it{r} between
\ili{Banda} and \ili{Geser} that led \citet[131ff]{Collins1986} to posit the “Proto-\ili{Banda}”
node of his East \ili{Central} Maluku that included \ili{Banda},
\ili{Geser} and \ili{Watubela}, with \ili{Watubela} reflecting a secondary development
of **r > \it{l}.
However, our \tsc{Eastern Islands} node of \tsc{Seram-Tanimbar-Bomberai}
includes additional data from \ili{Bati},
\ili{Kesui}, and \ili{Kowiai} (\srf{11-03}) data which were not available
and/or not examined by \citet{Collins1986}.
It is the behaviour of the apicals in \ili{Bati} (closely related to
\ili{Geser}-\ili{Gorom} with 80{\%} lexical similarity) that prevent us from
positing the merger of *d/*z/*R > **r for the whole \tsc{Eastern
Islands} node of \tsc{Seram-Tanimbar-Bomberai}.
\ili{Bati} merges **d/*l > \it{l}/{\#}{\gap} (word initial),
and **d/*l > \it{l},
{\0} /V{\gap}V (medial, intervocalic),
but has *R > \it{l} in all positions with no instances of medial
*R > {\0} (see \trf{11-16}).
This means we cannot posit a merger of **d/*R for \ili{Bati},
and by extension we cannot posit this merger for Proto-Eastern Islands.
The additional data from \ili{Bati} thus contravenes Collins’ claims
for “Proto-\ili{Banda}” based on the merger of *d/*z/*R > **r
and means we cannot use this merger to group \ili{Banda} with \tsc{Eastern Islands}.

We note that the \tsc{\ili{Teor}-Kur} node of \tsc{Seram-Tanimbar-Bomberai}
shares the merger of *d/*z/*R > **r with \ili{Banda},
as summarised in \trf{12-03}.
However, the considerable collective differences in the historical
phonology, the morphosyntax, and the lexicon lead us to conclude that this
merger shared by \ili{Banda},
\ili{Geser}, and \tsc{\ili{Teor}-Kur} is more likely to be due to parallel
development rather than to shared inheritance.

\begin{table}\tcp{PMP apicals in \ili{Banda} and \tsc{Seram-Tanimbar-Bomberai}}{12-03}
	\begin{threeparttable}\st{6pt}
	\begin{tabular}{lllllll}\lsptoprule
PMP &{\hr}*d&{\hr}*z&{\hr}*R&{\hr}*l&{\hr}*j\\	\midrule
\tsc{Seram-Tanimbar-Bomberai}&**d{\cbl}&**d\su{†}{\cbl}&**r{\cor}&**l{\cdr}&**j&\\
\hspace{5mm}\tsc{East Seram}&**l{\cdr}&**l{\cdr}&**r{\cor}&**l{\cdr}&**l{\cdr}&\\
\hspace{5mm}\tsc{Eastern Islands}&**d{\cbl}&**d{\cbl}&**r{\cor}&**l{\cdr}&**j&\\
\hspace{5mm}\tsc{Tanimbar-Bomberai}&**d{\cbl}&**d{\cbl}&**r{\cor}&**l{\cdr}&**j,&**r{\cor}\\
\hspace{5mm}\tsc{\ili{Teor}-Kur}&**r{\cor}&**r{\cor}&**r{\cor}&**l{\cdr}&**j&\\
\ili{Banda}&\hp{**}r{\cor}&\hp{**}r{\cor}&\hp{**}r{\cor}&\hp{**}l{\cdr}&\hp{**}l{\cdr}&\\
		\lspbottomrule
	\end{tabular}
		\begin{tablenotes}\tnsize
			\item[†] {While \tsc{Seram-Tanimbar-Bomberai} does have initial *d/*z > **d,
								the data are a little ambiguous for medial *d/*z > **d.
								See \srf{11-01-01} for discussion.}
		\end{tablenotes}
	\end{threeparttable}
\end{table}
